\begin{frame}{핵심 아이디어: 기댓값을 ``평균''으로 치환한다}
  가치함수의 정의:
  \[
  V^\pi(s) = \mathbb{E}_\pi[\,G_t \mid s_t = s\,]
  \]
  ``이 상태에서 시작해 정책을 따랐을 때 받을 리턴의
  \textbf{기댓값}.'' 기댓값을 정확히 계산하려면 전이확률이
  필요하지만 ---
  \vskip0.6em
  \begin{block}{통계학의 기본 원리 (대수의 법칙)}
    기댓값은 \textbf{많은 샘플의 평균}으로도 근사할 수 있다.
  \end{block}
  \vskip0.6em
  정책 $\pi$ 를 따라 에피소드를 여러 번 플레이하고, 상태 $s$ 를
  방문했을 때 이후 \textbf{실제로 받은 리턴}을 모아 평균 내면
  그것이 곧 $V^\pi(s)$ 의 추정치.
  \vskip0.4em
  \kq{$\mathbb{E}$ (확률분포를 ``알 때''의 연산)를 ``실제로 해본
  다음 평균 내는 것''(분포를 ``모를 때''의 연산)으로 치환해
  버린 것 --- 그것이 MC.}
\end{frame}
